\paragraph{Background on Tokenization.}
Tokenization is a simple way of mapping continuous representations
to discrete units. But why do we want discrete units in the first
place? Autoregressive models, such as GPT, are commonly trained
to predict discrete tokens using a cross-entropy loss. This makes
discrete representations a natural choice for autoregressive
action generation.

Traditionally, vector quantization (VQ)~\cite{vq} has been used
to obtain discrete representations. VQ maps continuous latents
to entries in a learned codebook. However, recent approaches
such as finite scalar quantization (FSQ)~\cite{fsq} provide
a simpler alternative. Rather than learning a codebook,
FSQ independently quantizes each latent dimension into a
finite set of values.

This induces a grid-like structure in the latent space,
providing a natural notion of distance and neighborhood
between discrete representations. We argue that this
structure provides a useful inductive bias for
autoregressive action generation.

Overall, the quantization scheme is simple and effective.